\documentclass[12pt,a4paper]{article}
\usepackage[spanish,es-tabla]{babel}
\decimalpoint
\usepackage{array}
\newcolumntype{C}[1]{>{\centering\arraybackslash}m{#1}}
\usepackage{bbm}
\usepackage{tikz}
\usetikzlibrary{fit,positioning}
\usepackage{pgfplots}
\pgfplotsset{compat=1.18}
\usepackage{amsmath}
\usepackage[utf8]{inputenc}
\usepackage{graphicx}
\usepackage{subcaption}
\usepackage{cancel}
\usepackage{multicol}
\usepackage{wrapfig}
\usepackage[T1]{fontenc}
\usepackage{url}
\usepackage{graphicx}
\usepackage{gensymb}
\usepackage{booktabs}
\usepackage{amssymb, amsmath}
\usepackage{wrapfig}
\parskip=1 mm
\oddsidemargin -1.5 cm
\headsep -2 cm
\textwidth=17.8cm
\textheight=25.5cm
\begin{document}
\begin{center}
    {\LARGE \textbf{Introducción a la Mecánica Cuántica}} \\[0.4cm]

    {\large Facultad de Ciencias, UNAM} \\[0.6cm]

    \large \textbf{Tarea 1} \\[0.2cm]

    \begin{tabular}{l}
        \small Autor: Fernando Naed Ortega Montero\\
        \small Fecha: \today
    \end{tabular}
\end{center}



\section*{Preguntas Conceptuales.}

\begin{enumerate}

\item Elija uno de los fenómenos estudiados (efecto fotoeléctrico, efecto Compton o digracción de electrones) y explique cómo determinaria experimentalmente la constante de Plank. Indique qué magnitudes necesita medir y como las utilizaría para obtener h.


\item ¿Qué evidencias experimentales disponibles en la época de Thomson llevaron a proponer un modelo en el que la carga positiva estuviera distribuida uniformemente los électrones se encontraran incrustados en ella?

Thomson realiza un experimento con los rayos catódicos, que al generar una diferencia de potencial entre los electrodos produce una linea de rayos catódicos.

Thomson observa que al aplicar campos magneticos o electricos se produce una desviación de los rayos catódicos, la cual coincide con que estos esten formados de partículas de con de carga electrica negativa.

Estas mediciones conducen a una proporción $\frac{e}{m_e} \approx 1.76 \times 10^{11} \frac{C}{kg}$, esta misma proporción fue medida con iones de hidrógeno $\frac{e}{m_p} \approx 10^8 \frac{C}{kg}$ lo cual es mucho menor que la proporción medida para los rayos catódicos, esto implica que $m_e << m_p$, así mismo observó que las propiedades de los rayos catódicos eran independientes de el gas dentro del tubo y del metal usado en los cátodos.

Su deducción fue que las particulas de las cuales estan echos los rayos catodicos son fundamentales en la compocición de la materia, y tienen mucha menos masa que los iones, por lo que tienen que formar parte de los átomos, pero el echo de que los atomos sean neutros y esta nueva particula que los compone tenga carga negativa, implica que el átomo esta formado de una carga positiva que neutralice la carga negativa.







\end{enumerate}


\end{document}